\section{Part II: Choosing an Agenda and Building a Research Portfolio}

\subsection{9. A research question must be falsifiable enough to guide allocation}

A good project charter states the current evidence, theoretical mechanism, competing explanations, measurable claim, appropriate baselines, data requirements, expected uncertainties, and possible results. A research question need not fit a yes/no experiment, but it must be specific enough that different findings would change beliefs. Set criteria for whether another experiment is warranted before the team becomes attached to a positive story.

For example, a new method that purportedly improves retrieval evaluation should specify evaluation datasets, leakage controls, ground truth quality, variance estimates, out-of-distribution tests, costs, and relevant existing methods. The result may be ``helps only on a narrow class of tasks,'' which is useful knowledge if reported honestly.

\subsection{10. Scientific merit, mission merit, and business merit are distinct}

Score proposed projects qualitatively against (a) \textbf{importance} of the question, (b) \textbf{plausibility} of producing reliable evidence, (c) \textbf{fit} with scientific talent and resources, (d) \textbf{value of learning}, (e) \textbf{downstream usefulness}, (f) \textbf{funding quality and restrictions}, (g) \textbf{ethical and security risk}, and (h) \textbf{opportunity cost}. Avoid summing arbitrary ordinal scores into a spurious decimal. Some projects fail an essential gate even if the sum is high.

A worthy mathematical question might have enormous scientific merit but low feasibility for a small applied institute. A funded client study may have good commercial merit but little publishable novelty. Label each honestly.

\subsection{11. Construct a portfolio by uncertainty type}

Different projects fail for different reasons. \textbf{Scientific risk:} the phenomenon may not exist. \textbf{Technical risk:} it exists but cannot be implemented economically. \textbf{Access risk:} data, equipment, participants, or permissions are unavailable. \textbf{Regulatory risk:} approval may not be granted. \textbf{Market risk:} no 1 adopts a useful invention. \textbf{Financing risk:} money ends before evidence arrives. Correlation matters: 5 AI experiments exposed to the same foundation-model shift are not 5 independent bets.

A balanced early portfolio could contain a few low-cost, high-learning applied inquiries and 1 more speculative option; a mature portfolio may include protected basic research and translation. Use position sizes and stop conditions rather than a rigid percentage rule.

\subsection{12. The exploration-exploitation trade-off}

Exploitation improves known activities; exploration tests uncertain possibilities. Both compete for the same cash and talent. A founder must protect the proven cash engine from a speculative research budget and protect research from being continually interrupted by routine paid work. The balancing mechanism is not motivational exhortation; it is \textbf{ring-fenced time, funds, and authority}, with explicit exception rules.

If a researcher spends every Friday on client emergencies, nominal ``1 day per week of research'' becomes zero. Conversely, a commercial division should not be forced to subsidize scientific work that management has never approved or measured.

\subsection{13. Staging scientific investment}

Use \textbf{discovery}, \textbf{replication}, \textbf{demonstration}, \textbf{translation}, and \textbf{deployment} gates. Discovery establishes whether a signal exists. Replication checks robustness. Demonstration asks whether the result operates under realistic conditions. Translation addresses engineering, IP, safety and adoption requirements. Deployment measures sustained value and consequences.

The gates need not be linear-fundamental research may never deploy, and a project may return from demonstration to discovery. Their value is preventing a striking demo from becoming a prematurely priced commercial asset.

\subsection{14. A project that should stop-and why stopping is a success}

[Model] A team spends ₹2 lakh to compare an experimental grading intervention with current expert feedback. The measured improvement on a pre-registered transfer assessment is close to zero, with uncertainty that excludes the minimum improvement needed to justify higher operating cost. The responsible response is to retire that method for this use case, document the negative result, and retain the evaluation procedure as an asset.

The team has created value by avoiding a ₹20-lakh platform build. A culture that insists ``we must launch something'' converts good research into bad capital allocation.

\subsection{15. Opportunity cost includes attention and institutional identity}

Every research initiative consumes the best people's focused attention, not just their billed hours. A new founder with 15-25 hours a week cannot simultaneously direct a serious lab, teach weekly, sell enterprise contracts, recruit faculty, and oversee compliance. Approving a project without a capable owner creates an invisible liability.

A useful investment memo lists (1) the most attractive alternative use of team time and (2) the institutional identity that the project reinforces. A project that earns a fee but forces the institute to become a confidential systems-integration shop should be assessed differently from 1 that strengthens a repeatable research-and-teaching core.

\subsection{16. Models for high-risk research portfolios}

For translation projects with plausible cash flows, scenario-weighted values and staged NPV can be useful. For basic science, economic capture may be too distant and uncertain; use budget constraints and externally reviewed scientific merit instead. Never force all projects into a ``risk-adjusted IRR'' fabricated from invented probabilities.

1 practical portfolio is \textbf{core evidence} (replicable work with near-term use), \textbf{developmental options} (promising methods with unresolved generalization), and \textbf{frontier inquiry} (important but highly uncertain science). Decide maximum loss and review cadence for each. Reserve the right to fund an important public-interest investigation with no plausible IP monetization-provided its subsidy is explicit and affordable.
